\section{Results and Discussion}
\subsection{Road Cost Agreement and Search Effort}
All 180 matched road queries return equal costs under A*, Dijkstra, and the frozen offline reference within $10^{-6}$ m. The graph validator reports zero heuristic edge violations and scale 1.0. Separately retained edge cases cover start equals goal, unreachable targets, directionality, and edge-midpoint snapping. These checks support correctness on the tested graph; agreement alone is not a proof for every possible adapter.

\begin{table}[H]
\caption{Frozen road results; 60 pairs per stratum. Expansion and time columns are marginal medians; reduction and ratio are medians of paired statistics.}\label{tab:road}
\centering\begin{tabular}{lrrrr}\toprule
Stratum & Expanded A*/D & Reduction & Time A*/D ($\mu$s) & Ratio\\\midrule
Short &101 / 505&78.6\%&44.95 / 128.09&0.324\\
Medium &301 / 1,592&77.4\%&282.95 / 1,117.24&0.287\\
Long &714.5 / 2,588.5&68.6\%&771.46 / 1,883.19&0.449\\\bottomrule
\end{tabular}
\end{table}

\begin{figure}[t]
\centering\includegraphics[width=\textwidth]{road-effort.pdf}
\caption{Matched road search effort and search-call timing. Points below the diagonal favor A*. Each stratum has 60 queries; all 180 costs agree. Timing excludes graph loading, snapping, API, and rendering.}
\label{fig:road-effort}
\end{figure}

Table~\ref{tab:road} and Fig.~\ref{fig:road-effort} show fewer expansions in every A* run, with median paired reduction 76.0\%. A* is faster in 173 cases, Dijkstra in seven; the median time ratio is 0.338, range 0.061--1.435, with pair 167 the largest regression. The heuristic improves descriptive effort and time on this network without winning every timing pair. Paired ratios are not ratios of marginal medians. This known-graph comparison controls representation more tightly than the robot study.

\subsection{Robot Local Geometry and Path Length}
Both local planners find paths for all 21 eligible requests, with correct endpoints and primary interior-only validity in all 21. These are local planning successes, not 21 successful autonomous missions. Of the ten prospective source episodes, nine reach the goal and one reaches its iteration cap; its eligible states remain. Figure~\ref{fig:observed} shows an observed-space request and the two returned paths. No hidden polygons are supplied to either local planner.

\begin{figure}[t]
\centering\includegraphics[width=\textwidth]{robot-observed.pdf}
\caption{Representative frozen dead-end request. Visible blocker fragments and approximate display sight footprints show recorded knowledge; ASP and scan-center A* use the same position and target. Sight footprints are display approximations; benchmark link admission uses analytic whole-segment certification. Hidden obstacles are omitted.}
\label{fig:observed}
\end{figure}

\begin{table}[H]
\caption{Local robot outcomes by cohort. ``Shorter'' counts follow the primary point convention. Two prospective cases have equal lengths. Delta is the median paired A* minus ASP length in world units.}\label{tab:robot}
\centering\begin{tabular}{lrrrr}\toprule
Cohort &Pairs&ASP shorter&A* shorter&Delta\\\midrule
Initial &6&5&1&$+7.403$\\
Prior held-out &1&1&0&$+1.654$\\
Prospective &14&3&9&$-2.492$\\\midrule
Combined &21&9&10&$0.000$\\\bottomrule
\end{tabular}
\end{table}

ASP is shorter in nine requests, A* in ten, and two are equal. The paired delta median is zero, mean $+2.981$, range $[-10.707,+44.266]$ world units. Cohort composition matters: five of six initial cases favor ASP, nine of fourteen prospective cases favor A*. One held-out eligible request cannot establish generalization. Figure~\ref{fig:robot-length} retains all pairs.

\begin{figure}[t]
\centering\includegraphics[width=\textwidth]{robot-length.pdf}
\caption{All 21 eligible local requests, grouped by frozen cohort. Positive A* minus ASP differences favor ASP. Filled orange markers indicate ASP boundary contact under the strict convention; these paths remain primary interior-valid.}
\label{fig:robot-length}
\end{figure}

Strict no-contact validity differs: ASP passes 15 of 21, and A* passes 21 of 21. All six boundary-touching ASP paths are shorter under the permissive point convention. In the largest difference, \source{P02-_map_deadend-01}, ASP is 15.734 units versus A* 60.000, but ASP touches a boundary. Among the 15 jointly strict-valid pairs, ASP is shorter in three, A* in ten, and two are equal; the median paired difference is $-2.584$. This secondary analysis changes the feasible interpretation and must not replace the prespecified primary result. It illustrates how strict links through sparse centers and geometric refinement produce different paths.

Radius-0.5 clearance passes only nine ASP and eleven A* paths, precluding disk-safety claims. Significant-turn medians are two and one respectively; paired all-turn delta median is $-2$. Tiny ASP bends can inflate all-turn counts without corresponding physical steering. ASP vertices can lie inside bundles absent from the sparse A* graph, so mixed lengths do not contradict graph-optimal search.

\begin{figure}[t]
\centering\includegraphics[width=\textwidth]{robot-validity.pdf}
\caption{Validity outcomes for the same 21 requests. Primary point validity allows boundary tangency; strict validity forbids it. Disk clearance is only a post-plan sensitivity check, not a guarantee of executed finite-robot motion.}
\label{fig:validity}
\end{figure}

\subsection{Robot Computation Components}
Figure~\ref{fig:runtime} separates the measured components. Across state medians, ASP core time is 62.655 ms, critical-line construction 0.319 ms, and optimization 62.400 ms. A* link-build time is 42.102 ms, binding/search 0.0263 ms, and the median per-state build-plus-call sum 42.142 ms. The small C++ call is a search on already constructed sparse graphs, while Python geometric certification can dominate its surrounding pipeline.

\begin{figure}[t]
\centering\includegraphics[width=\textwidth]{robot-runtime.pdf}
\caption{Per-state median local computation times (three warm-ups, 31 calls). A logarithmic axis exposes component scale. ASP core excludes skeleton generation; A* build includes observed-link certification. Components and languages differ, so isolated search ratios are not algorithm-only speedups.}
\label{fig:runtime}
\end{figure}

Dividing ASP core time by the C++ call would misrepresent differing languages and preparation. Even build-plus-call and ASP core omit different upstream work. The justified conclusion is that construction matters and C++ search is small for these requests, not that an intrinsic algorithm speedup was established.

\subsection{Supporting Online Behavior and Recovery}
The legacy and radar policies both reach the goal in 20 paired episodes. Radar travels less in 18 and more in two, with median paired delta $-95.419$ world units. Its largest regression is $+70.772$ in Complex Irregular Labyrinth at radius 7. This supports retaining radar ranking as a useful non-dominating policy. It does not prove an optimal exploration order or justify exhausting low-score frontiers. Twelve sensor-radius episodes all reach their goals, but the radius-5 and radius-7 distance medians are 92.909 and 113.278 units; greater range does not monotonically shorten travel in this sample.

Cached and full frontier construction preserve complete playback, recovery outcomes, status, and distance in all 20 pairs. Their median frontier-processing times are 639.166 and 46.723 ms respectively. These are accumulated component times per episode, not end-to-end timings. Cache validity depends on invalidation when newly observed free geometry or walls change the frontier mask. The paired equivalence is concrete evidence for the tested indoor and maze configurations, not a universal assertion about arbitrary geometric updates.

\begin{figure}[t]
\centering\includegraphics[width=\textwidth]{supporting.pdf}
\caption{Supporting evidence from our online simulator. Left: paired frontier processing with equivalent playback in 20 cases. Right: four observed indoor branch returns, each certified and marked nonfallback in the frozen trace. These are separate from the local authors-versus-A* comparison.}
\label{fig:support}
\end{figure}

Four certified nonfallback indoor returns demonstrate genuine C++ A* recovery. Verified reverse-entry fallback is distinct. Certification concerns executed length to an observed ancestor; successful unlabelled exits, unreachable fixtures, and budget termination require separate labels. Continuous no-frontier status establishes candidate exhaustion, not unknown-world unreachability.
